\section{Introduction}
\label{sec:intro}

Which convex polytope with $n$ vertices, inscribed in the unit sphere
$\mathbb{S}^{d-1}$ of $\mathbb{R}^{d}$, has the largest volume? Fejes T\'{o}th
raised this question in three dimensions \cite{fejes-toth-1964}, and it is
still open there for most $n$ (see \cite[Section~1]{horvath-langi-2016}).
There, besides the classical cases $n = 4, 6,
12$ (see \cite[Section~1.2]{hoehner-ledford-2026}), Berman and Hanes
determined the optimal polyhedra with at most eight vertices
\cite{berman-hanes-1970}, and Hoehner and Ledford recently those with nine
\cite{hoehner-ledford-2026}. For more vertices only numerical results are
known: Mutoh reported a computer search for up to $30$ vertices
\cite{mutoh-2003}, and Hardin, Sloane and Smith tabulate putatively optimal
configurations for up to $130$ \cite{hardin-sloane-smith}.

In higher dimensions even less is known. Horv\'{a}th and L\'{a}ngi extended
the necessary condition for optimality of Berman and Hanes to every dimension,
and with it showed that optimal polytopes are simplicial
\cite[Lemma~1]{horvath-langi-2016}. They determined the optimal polytopes with
$n = d + 1$ vertices, the regular simplices
\cite[Corollary~1]{horvath-langi-2016}; with $n = d + 2$, the convex
hulls of two regular simplices of dimensions $\lfloor d/2\rfloor$ and
$\lceil d/2\rceil$ in orthogonal subspaces
\cite[Theorem~2]{horvath-langi-2016}; and, for odd $d$, with $n = d + 3$
\cite[Theorem~3]{horvath-langi-2016}. For even $d$ their argument for
$n = d + 3$ excludes cyclic polytopes, and they asked whether an optimal
polytope can be cyclic \cite[Problem~2]{horvath-langi-2016}.

In dimension four these results settle $n = 5$ and $n = 6$. In earlier work
we answered the question of Horv\'{a}th and L\'{a}ngi for $d = 4$, which
settles $n = 7$, and showed that the regular cross-polytope is optimal for
$n = 8$ \cite{medgyes-maxvol-d4}. We then settled $n = 9$
\cite{medgyes-maxvol-d4-n9}. The optima for $n = 6$ to $9$ share a structure:
each is the convex hull of two regular polygons with $m$ and $l$ vertices,
inscribed in $\mathbb{S}^{3}$ in orthogonal $2$-planes through the origin,
with $(m, l) = (3, 3)$, $(3, 4)$, $(4, 4)$ and $(4, 5)$. The method bounds
the volume on each combinatorial type of simplicial polytope by a
sum-of-squares certificate \cite{parrilo-2003}, an algebraic identity checked
in exact rational arithmetic, and treats the type of the optimum separately.
To our knowledge, no case with $d = 4$ and $n \geq 10$ has been settled.

This paper treats $n = 10$, where the pattern suggests a pair of regular
pentagons. Two things change from nine vertices. The number of simplicial
types grows from $1142$ \cite{altshuler-bokowski-steinberg-1980} to $162004$
\cite[Theorem~25(i) and Table~3]{firsching-2017}. And the margin is small:
to exclude the types other than that of the pair, we need a bound less than
$2 \cdot 10^{-4}$ below its volume. At this margin the certificates of
\cite{medgyes-maxvol-d4}, which bound the signed volume on every
configuration of points on the sphere, convex or not, fail on one type, the
cyclic polytope $C_{4}(10)$.

We prove that the pair of regular pentagons is, up to orthogonal
transformations, the unique optimum (Theorem~\ref{thm:main}). On the type of
the pair, the join of two $5$-cycles, the sharp bound follows from
\cite[Corollary~5.2]{medgyes-maxvol-d4-n9}. For $C_{4}(10)$ we add
nonnegative multiples of facet polynomials, which are positive on every
realisation of a type, and obtain a \emph{convexity certificate}
\cite[Section~5]{parrilo-2003}: it bounds only the volume of the polytopes of
the type, which is all the problem requires.
